\documentclass{situationist-preprint}

\title{A Détournement of the Preceding Manuscript}
\msid{POL-24-b}
\msfamily{POL-24}
\msdiscipline{Political Philosophy}
\mstype{Companion manuscript}
\mskeywords{détournement; recuperation; spectacle; Debord; Wolman; quotation}

\begin{document}
\maketitle

\begin{abstract}
POL-24-a argues that the Situationists would regard the present collection as an instance of the spectacle, a representation of thought that takes the place of thinking, and it leaves to the present manuscript the question whether a détournement can be released in the same series as its target. We take POL-24-a's own sentences and redeploy them against it, in the sense of Debord and Wolman; they show that the critique of the collection is itself received as an item of the collection. A détournement can be released in the same series as its target, but it is then a representation of détournement. This capture is an instance of the conclusion of POL-24, not a refutation of it.
\end{abstract}

\section{Détournement}

\citet{DebordWolman1956} describe détournement as the reuse of pre-existing elements in a new ensemble that gives them a new meaning. Détournement, they hold, is weaker the closer it comes to a rational reply. In \emph{The Society of the Spectacle}, \citet[theses~206--209]{Debord1967} restates the method: détournement is the opposite of quotation, which invokes an authority already falsified by having been torn from its context.

The present manuscript is held to a rule that détournement does not obey: its quotations must be exact and cited. Each diversion below therefore sets a sentence of POL-24-a, quoted exactly, beside the sentence it becomes when redirected at POL-24-a itself.

\section{Three diversions}

\emph{D1.} POL-24-a: \enquote{a reader who learns that a theorem has been verified, without reading the list of its axioms, receives a further representation, of rigour in place of rigour.} Diverted: a reader who learns that the collection has been criticized, without reading the criticism, receives a further representation, of critique in place of critique. The reader of the catalogue line of POL-24 learns that the collection has examined itself, and what it found, without the examination taking place for her.

\emph{D2.} POL-24-a: \enquote{That relation is not one of joint inquiry, in which each party may change the other's mind, but one of release and reception.} Diverted: the relation between POL-24-a and the present manuscript is not one of joint inquiry, but one of release and reception. POL-24-a is fixed, the present manuscript is bound to its conclusion, and its title, type and length were assigned in advance in the catalogue. A détournement listed in a catalogue, with a target length, before it is written, is a commission.

\emph{D3.} POL-24-a: \enquote{The one point at which the diagnosis does not apply is the point at which a reader stops receiving the collection and starts arguing with it.} Diverted: the point at which a manuscript of the collection starts arguing with it is a point at which the collection is received. The present manuscript argues with POL-24-a, as item POL-24-b, in the form of everything else the collection releases.

\section{Capture}

\begin{proposition}\label{p:capture}
A détournement can be released in the same series as its target; so released, it is spectacular in the sense of Definition~1.1 of POL-24-a.
\end{proposition}

\begin{proof}
The present manuscript is such a release. As to the marks: (M1) For a reader of the catalogue, the entry \enquote{A Détournement of the Preceding Manuscript} takes the place of the diversion, by D1. (M2) The relation between the manuscript and its readers, and between it and its target, is mediated by the series' metadata, by D2 and D3. (M3) It is produced and circulated under the conditions that POL-24-a shows the whole collection to satisfy.
\end{proof}

The constraint of exact quotation already marks the change: by Debord's account in thesis~208, a diversion that must quote and cite its source has become the quotation it was meant to oppose.

\begin{proposition}\label{p:confirms}
The capture of the present manuscript confirms the conclusion of POL-24.
\end{proposition}

\begin{proof}
POL-24-a's reply to its second objection holds: \enquote{The spectacle does not suppress its critique; it represents it, and the representation of critique is one of its most effective forms}. Had the present manuscript escaped the form of the series, that reply would have a counterexample; by Proposition~\ref{p:capture} it did not, and it is a further instance of what the reply describes, and so of the conclusion of POL-24.
\end{proof}

\begin{objection}
If every critique of the collection, even a failed one, confirms that it is spectacular, the conclusion of POL-24 is irrefutable, and says nothing.
\end{objection}

\begin{reply}
It is refutable, and POL-24-a's first reply states how: for a reader who checks an argument or follows a reference to its source, the first mark fails. What refutes the conclusion is a way of reading, not a text. No manuscript of the series can supply it, since every manuscript is received as the others are. A reader who opens POL-24-a to check whether the diversions above are faithful to it supplies it, by an activity that the present manuscript can occasion and cannot perform.
\end{reply}

\section{Conclusion}

The détournement of POL-24-a succeeds as an argument and fails as a détournement: released in the same series as its target, it becomes a representation of one, as POL-24-a predicted. The collection remains an instance of the spectacle: a representation of thought that takes the place of thinking.

\begin{thebibliography}{9}\small
\bibitem[Debord(1967)]{Debord1967} Debord, G. (1967). \emph{La Société du spectacle}. Buchet-Chastel. Trans. \emph{The Society of the Spectacle}, Zone Books, 1994.
\bibitem[Debord and Wolman(1956)]{DebordWolman1956} Debord, G. and Wolman, G.~J. (1956). Mode d'emploi du détournement [A user's guide to détournement]. \emph{Les Lèvres nues}, 8.
\end{thebibliography}

\end{document}
